
by a commutative \(S\)-group prescheme \(D(G)\), finite and flat over
\(S\), and the canonical morphism
\[
  G\longrightarrow D(D(G))
\]
is an isomorphism.  In particular:
\begin{enumerate}
\item[(i)] \(D(\boldsymbol\mu_n)\simeq\mathbb Z/n\mathbb Z\), and
  consequently \(D(\mathbb Z/n\mathbb Z)\simeq\boldsymbol\mu_n\);
\item[(ii)] if \(S\) has characteristic \(p>0\), then
  \(D(\boldsymbol\alpha_p)\simeq\boldsymbol\alpha_p\).
\end{enumerate}

\SGASection{3}{Unipotent groups acting on a vector space}
\SGARefTarget{sga3:xvii:section:3}

Recall (\SGARefLink{sga3:i:definition:4-6-1}{Exposé I, 4.6.1}) that if \(S\) is a prescheme and \(M\) is a
sheaf of \(\mathcal O_S\)-modules, \(\mathbf W(M)\) denotes the
\(S\)-functor
\[
  (\operatorname{Sch}/S)^\circ\longrightarrow(\operatorname{Ens})
\]
defined, for every \(S\)-prescheme \(S'\), by
\[
  \mathbf W(M)(S')
  =\Gamma\!\left(S',M\otimes_{\mathcal O_S}\mathcal O_{S'}\right).
\]

Recall also that if an \(S\)-group \(G\) acts on an \(S\)-functor
\(V\), the \(S\)-functor \(V^G\) of invariants of \(V\) under \(G\) is
the subfunctor of \(V\) whose \(S'\)-valued points, for \(S'\) over
\(S\), are those \(x\in V(S')\) such that \(x_{S''}\) is fixed under
\(G(S'')\) for every \(S''\) over \(S'\).

We then have the following lemma.

\SGAStatement{\SGAInlineRefTarget{sga3:xvii:lemma:3-1}Lemma 3.1}{Let \(S\) be a prescheme, and let \(G\) be an
\(S\)-group prescheme, affine over \(S\), defined by the
quasi-coherent \(\mathcal O_S\)-algebra \(A\).  Suppose that \(G\) acts
on a quasi-coherent sheaf \(M\) of \(\mathcal O_S\)-modules.  Let
\[
  \mu:M\longrightarrow A\otimes_{\mathcal O_S}M
\]
be the comorphism defining the action of \(G\) on \(M\) (\SGARefLink{sga3:i:proposition:4-7-2}{Exposé I,
4.7.2}), and let
\[
  \nu:M\longrightarrow A\otimes_{\mathcal O_S}M
\]
be the morphism \(x\mapsto\mu(x)-1\otimes x\).  Then:
\begin{enumerate}
\item[\SGAInlineRefTarget{sga3:xvii:lemma:3-1:item:i}(i)] \(\underline M^{\,G}(S)=\Gamma\operatorname{Ker}(\nu)\);
\item[\SGAInlineRefTarget{sga3:xvii:lemma:3-1:item:ii}(ii)] if \(S\) is the spectrum of a field \(k\), then
  \(\underline M^{\,G}\) is of the form \(\mathbf W(N)\), where \(N\)
  is a vector subspace of \(M\);
\item[\SGAInlineRefTarget{sga3:xvii:lemma:3-1:item:iii}(iii)] if \(S\) is the spectrum of a field \(k\), every element
  \(x\) of \(M\) is contained in a finite-dimensional \(k\)-vector
  subspace of \(M\) stable under the action of \(G\).
\end{enumerate}}
\SGARefTarget{sga3:xvii:lemma:3:1}

\SGAProof{Proof}
Assertion \SGARefLink{sga3:xvii:lemma:3-1:item:iii}{(iii)} is recalled here for reference and was already proved
in Exposé VI\textsubscript{B}, \SGARefLink{cor:sga3-VIB-11.2}{11.2}.

For \SGARefLink{sga3:xvii:lemma:3-1:item:i}{(i)}, it is clear that \(\underline M^{\,G}(S)\) contains
\(\Gamma\operatorname{Ker}(\nu)\).  To prove the reverse inclusion, we
may suppose that \(S\) is affine, with ring \(B\).  Let
\(m\in\underline M^{\,G}(S)\).  For every \(B\)-algebra \(B'\) and
every
\[
  u\in\operatorname{Hom}_{B\text{-alg}}(A,B')
\]
(where \(u\) corresponds to an element of
\(G(\operatorname{Spec}B')\)), we have
\[
  (u\otimes1_M)\nu(m)=0
  \qquad\text{in }B'\otimes_BM.
\]
Taking \(B'=A\) and \(u\) to be the identity of \(A\), we obtain
\(\nu(m)=0\).

For (ii), let \(N\) be the kernel of \(\nu\), equal to
\(\underline M^{\,G}(k)\) by \SGARefLink{sga3:xvii:lemma:3-1:item:i}{(i)}.  Every \(k\)-prescheme \(S\) is flat
over \(k\); hence
\[
  \underline M^{\,G}(S)
  =\Gamma\operatorname{Ker}(\nu\otimes_k S)
  =\Gamma(N\otimes_k S).
\]
Thus \(\underline M^{\,G}\) is isomorphic to the functor
\(\mathbf W(N)\).

\SGAStatement{Proposition 3.2}{Let \(G\) be a unipotent algebraic
group, defined over a field \(k\), acting on a \(k\)-vector space
\(V\).  If \(V\ne0\), then \(V^G\ne0\).}
\SGARefTarget{sga3:xvii:proposition:3:2}
